\section{Pillar III --- The Analysis of Phase Transitions}
\label{sec:analysis}

\Cref{ch:introduction} stakes the thesis on a metaphor: convergence is a phase
transition, like water becoming steam. A metaphor earns its keep only if the
mathematics it points to is real. This pillar supplies that mathematics from first
principles: the calculus of continuity and discontinuity, the thermodynamic
classification of transitions, a derivation of the Clausius--Clapeyron relation
that shows the steam expansion ratio and the latent heat are \emph{one equation},
and a Landau model in which an order parameter jumps discontinuously as a control
parameter crosses a threshold.

\subsection{Limits, continuity, and discontinuity}
\label{sec:an-limits}

\begin{definition}[Limit and continuity]
A function $f\colon\mathbb R\to\mathbb R$ has \emph{limit} $\ell$ at $x_0$, written
$\lim_{x\to x_0}f(x)=\ell$, if for every $\varepsilon>0$ there is $\delta>0$ with
$0<|x-x_0|<\delta\Rightarrow|f(x)-\ell|<\varepsilon$. It is \emph{continuous} at
$x_0$ if $\lim_{x\to x_0}f(x)=f(x_0)$, and \emph{differentiable} there with
derivative $f'(x_0)=\lim_{h\to0}\frac{f(x_0+h)-f(x_0)}{h}$ when that limit exists.
\end{definition}

\begin{definition}[Jump discontinuity]
$f$ has a \emph{jump discontinuity} at $x_0$ if the one-sided limits
$f(x_0^-)=\lim_{x\uparrow x_0}f(x)$ and $f(x_0^+)=\lim_{x\downarrow x_0}f(x)$ exist
but differ; the \emph{jump} is $f(x_0^+)-f(x_0^-)\ne 0$.
\end{definition}

A jump discontinuity is the exact sense in which the thesis claims the order
parameter $\Phi$ of \cref{sec:fw-phase} is ``discontinuous'' in the control
$\sigma$: not that it changes quickly, but that its one-sided limits at $\sigma_c$
differ. This is what distinguishes a phase transition from rapid but smooth growth.

\subsection{Thermodynamic potentials and the Ehrenfest classification}
\label{sec:an-thermo}

\begin{definition}[Gibbs free energy]
For a substance at temperature $T$ and pressure $P$, the molar \emph{Gibbs free
energy} is $g=u-Ts+Pv$, with $u$ internal energy, $s$ molar entropy, $v$ molar
volume. Its differential is $\mathrm dg=-s\,\mathrm dT+v\,\mathrm dP$, so that
$s=-\bigl(\partial g/\partial T\bigr)_P$ and
$v=\bigl(\partial g/\partial P\bigr)_T$.
\end{definition}

\begin{definition}[Ehrenfest classification]
A phase transition is \emph{first-order} if the first derivatives of $g$
(equivalently $s$ and $v$) are discontinuous across it, and \emph{continuous}
(``second-order'') if the first derivatives are continuous but some second
derivative is not. A first-order transition therefore involves a jump $\Delta s$
in entropy and $\Delta v$ in volume.
\end{definition}

Water-to-steam is the prototypical first-order transition: at the boiling point
the molar volume jumps by the factor of ${\sim}1{,}600$--$1{,}700$ recorded in
\cref{sec:metaphor-intro} \citep{engtoolbox-steam,nationalboard-steam}, and the
entropy jumps by $\Delta s=L/T$ where $L$ is the latent heat
\citep{latentheat-openstax}. The thesis's claim is precisely that the
standardisation transition is \emph{first-order} in this sense---a jump in the
addressable surface, not a kink.

\subsection{The Clausius--Clapeyron relation, from first principles}
\label{sec:an-clausius}

\begin{theorem}[Clausius--Clapeyron]
\label{thm:clausius}
Along a coexistence curve $P(T)$ where two phases are in equilibrium, the slope
satisfies
\[
\frac{\mathrm dP}{\mathrm dT}=\frac{\Delta s}{\Delta v}=\frac{L}{T\,\Delta v},
\]
where $\Delta s=s_2-s_1$, $\Delta v=v_2-v_1$, and $L=T\,\Delta s$ is the latent
heat absorbed at constant $T$.
\end{theorem}
\begin{proof}
At coexistence the two phases have equal molar Gibbs free energy,
$g_1(T,P)=g_2(T,P)$. Moving along the coexistence curve preserves this equality,
so $\mathrm dg_1=\mathrm dg_2$. Substituting
$\mathrm dg_i=-s_i\,\mathrm dT+v_i\,\mathrm dP$ gives
$-s_1\,\mathrm dT+v_1\,\mathrm dP=-s_2\,\mathrm dT+v_2\,\mathrm dP$, hence
$(v_2-v_1)\,\mathrm dP=(s_2-s_1)\,\mathrm dT$ and
$\mathrm dP/\mathrm dT=\Delta s/\Delta v$. For a reversible change at constant
temperature the absorbed heat is $L=T\,\Delta s$, so $\Delta s=L/T$ and the second
equality follows.
\end{proof}

\begin{corollary}[Expansion and latent heat are one equation]
\label{cor:onesfact}
Rearranging \cref{thm:clausius}, $L=T\,\Delta v\,(\mathrm dP/\mathrm dT)$. For
vaporisation $v_{\text{vapour}}\gg v_{\text{liquid}}$, so $\Delta v\approx
v_{\text{vapour}}$ and the latent heat is, to leading order, proportional to the
volume expansion. The two numbers the thesis cites for water---an expansion of
${\sim}1{,}600$--$1{,}700\times$ and a latent heat of ${\sim}2{,}256\
\mathrm{kJ/kg}$---are not independent facts but two readings of
\cref{thm:clausius}.
\end{corollary}

\Cref{cor:onesfact} is the rigorous core of the metaphor. The ``latent heat of
standardisation'' (\cref{def:latent}) is not a loose analogy bolted onto a volume
figure: in the physical system the large expansion \emph{is} the large latent
heat, related by \cref{thm:clausius}. The thesis's two claims---that
standardisation costs a great deal (latent heat) and that it yields a great
expansion (volume)---are, under the metaphor, the same claim.

\subsection{Landau theory: a discontinuous jump from a smooth potential}
\label{sec:an-landau}

The Ehrenfest picture describes a transition; Landau theory \emph{generates} one
from a free energy that is itself smooth in an order parameter $\varphi\ge 0$.

\begin{construction}[Landau free energy with a cubic term]
\label{con:landau}
Let the (coarse-grained) free energy be
\[
F(\varphi)=a\,\varphi^{2}-b\,\varphi^{3}+c\,\varphi^{4},\qquad b,c>0,\ a\ge 0,
\]
where the control parameter (here, inverse standardisation, decreasing as
$\sigma$ rises) is encoded in $a$. The equilibrium order parameter is
$\varphi^{*}(a)=\arg\min_{\varphi\ge0}F(\varphi)$.
\end{construction}

\begin{proposition}[First-order jump]
\label{prop:landau-jump}
For \cref{con:landau}, the global minimiser $\varphi^{*}(a)$ is discontinuous at
$a_c=\dfrac{b^{2}}{4c}$: for $a>a_c$ the minimiser is $\varphi^{*}=0$, and as $a$
decreases through $a_c$ it jumps to $\varphi^{*}=\dfrac{b}{2c}>0$.
\end{proposition}
\begin{proof}
Stationarity: $F'(\varphi)=\varphi\,(2a-3b\varphi+4c\varphi^{2})=0$, so
$\varphi=0$ is always stationary, with other roots solving
$4c\varphi^{2}-3b\varphi+2a=0$. A nonzero minimiser becomes \emph{global} when its
free energy equals that at the origin, $F(\varphi)=F(0)=0$, i.e.\ (dividing the
quartic by $\varphi^{2}$) $a-b\varphi+c\varphi^{2}=0$. Solve the two conditions
simultaneously: from $a-b\varphi+c\varphi^{2}=0$ we get $a=b\varphi-c\varphi^{2}$;
substituting into $4c\varphi^{2}-3b\varphi+2a=0$ yields
$4c\varphi^{2}-3b\varphi+2b\varphi-2c\varphi^{2}=2c\varphi^{2}-b\varphi
=\varphi(2c\varphi-b)=0$, so $\varphi=b/(2c)$. Back-substituting,
$a_c=b\cdot\frac{b}{2c}-c\cdot\frac{b^{2}}{4c^{2}}
=\frac{b^{2}}{2c}-\frac{b^{2}}{4c}=\frac{b^{2}}{4c}$. For $a>a_c$ the origin is the
unique global minimum; for $a<a_c$ the branch $\varphi=b/(2c)+O(a_c-a)$ has strictly
lower free energy. Hence the global minimiser jumps from $0$ to $b/(2c)$ as $a$
crosses $a_c$ from above.
\end{proof}

\Cref{prop:landau-jump} is a complete, from-scratch instance of the phenomenon the
thesis needs: a free energy that depends \emph{smoothly} on its parameters
nonetheless yields an order parameter that is a \emph{discontinuous} function of
the control. The cubic term $-b\varphi^{3}$---the analogue of a network effect, in
which partial adoption is disproportionately rewarded---is what makes the
transition first-order rather than continuous.

\subsection{The standardisation transition}
\label{sec:an-standardisation}

Assembling the pillars: let $\varphi$ be the fraction of the language ecosystem
operating on the fused substrate, and let $a$ decrease as standardisation
$\sigma$ and adoption incentives rise. \Cref{prop:landau-jump} gives a critical
$\sigma_c$ (where $a=a_c$) at which $\varphi$ jumps from $0$ to a finite value:
the discontinuity asserted in \cref{prop:discontinuous}. The size of the resulting
addressable surface is governed not by $\varphi$ alone but by the algebraic gearing
of Pillar~I: by \cref{cor:superlinear}, once a fraction $\varphi$ of $n{\approx}$
domains, capabilities, and agents is interconnected, the reachable surface scales
as $\Theta((\varphi n)^{3})$ against $\Theta(\varphi n)$ integration cost. A
discontinuous jump in $\varphi$ (Pillar~III) multiplied by a cubic gearing in
reach (Pillar~I) is the formal content of the water-to-steam image: a threshold
crossed, and an order-of-magnitude expansion released.

\begin{remark}[On honest application]
Three cautions, consistent with \cref{prin:plateau} and \cref{sec:m-epistemics}.
First, $\Phi$ for a sociotechnical ecosystem is not measured with the precision of
a steam table; the model is explanatory, and any numeric projection of $\Phi$ is a
\forecast\ (\cref{ch:vision}). Second, \cref{prop:landau-jump} also admits the
case $a_c$ never being reached---a plateau that does not end in a transition.
Third, the analogy is structural, not literal: we claim the convergence shares the
\emph{order type} of a first-order transition (a jump), established by a
mechanism (threshold $\times$ gearing) that is itself rigorous, not that
information obeys thermodynamics. With those cautions, the metaphor of
\cref{ch:introduction} is now a model with stated mechanism, derived from first
principles. Contemporary work reading learning itself through phase transitions
and renormalisation (surveyed in \cref{sec:recent}) is the empirical company this
model keeps.
\end{remark}
